\section{Read the envelope before interpreting its payload}
The eight signature bytes are hexadecimal \texttt{89 50 4E 47 0D 0A 1A 0A}. After them,
a chunk begins with a big-endian length. The length counts only payload bytes. For
IHDR, that length is thirteen, so the complete chunk occupies 25 bytes. The next
chunk begins at file offset 33: eight signature bytes plus those 25 bytes.

\Needspace{10\baselineskip}
\begin{longtable}{p{24mm}p{35mm}p{91mm}}
\toprule Part & Byte count & Interpretation \\\midrule
Length & 4 & Unsigned, big-endian payload size \\
Type & 4 & Four ASCII letters; included in CRC \\
Payload & Declared length & Chunk-specific data; included in CRC \\
CRC & 4 & Big-endian stored CRC of type and payload \\\bottomrule
\end{longtable}
A big-endian field is a byte convention, independent of the simulator's memory byte
order. Assemble its numerical value from individual bytes. For example,
\texttt{00 00 02 00} means 512 in either target byte order. Validate that value as a
count before using it as a dimension. A machine word load from the file may be both
unaligned and interpreted in the wrong byte order.

The type's first letter distinguishes critical from ancillary chunks through its
ASCII case bit. Test this property only after checking that the bytes are letters.
A matching name alone does not establish legal placement: a second IHDR still fails.
Maintain explicit state for header seen, palette seen, IDAT not started/open/closed,
and IEND reached. Test the transition before processing a known chunk's payload.
